\chapter{抽象代数基础}
尽管本书侧重以几何直观作为理解线性代数的起点，但仍需借助准确的语言和严谨的符号加以表达。同时，线性代数中的部分概念也建立在基本代数结构之上。因此，本章将从集合与映射出发，简要介绍群与域，并讨论与后文相关的基本性质。这些内容既是对初等教育阶段所学数学知识的抽象与整理，也将帮助读者建立统一的符号语言，为后续学习做好准备。

\section{集合与映射}

在初等数学中，我们早已接触过集合（set）的概念，例如自然数集、实数集等等。那么，集合究竟是什么呢？在朴素的观念中，我们通常认为集合就是``符合任意特定性质的一堆东西''。然而，这一观念会导致悖论的出现，为此，公理化集合论（Zermelo–Fraenkel set theory, ZFC）引入了严格的公理体系来规定集合必须满足的条件。具体而言，一个合法的集合应当满足以下公理：
\begin{itemize}
    \item \textbf{外延公理}：两个集合如果含有完全相同的元素，那么它们就是同一个集合。
    \item \textbf{空集公理}：存在一个不含任何元素的集合，记作 \(\varnothing\)。
    \item \textbf{配对公理}：任意两个集合 \(A,B\) 可以放在一起，组成一个新的集合 \(\{A,B\}\)。
    \item \textbf{并集公理}：如果一个集合的元素本身也是集合，那么可以把这些集合中的元素合并起来，组成一个新的集合。
    \item \textbf{幂集公理}：任意集合 \(A\) 的所有子集可以组成一个新的集合，记作 \(\mathcal P(A)\)。
    \item \textbf{分离公理}：可以从一个已经存在的集合中，挑出满足某个条件的元素，组成一个新的集合。
    \item \textbf{替换公理}：如果按照某个确定规则，把一个集合中的每个元素都变成另一个对象，那么这些结果也可以组成一个新的集合。
    \item \textbf{无穷公理}：存在无限集合，因此可以构造自然数集合。
    \item \textbf{正则公理}：集合不能包含自己，也不能出现一层一层无穷向下嵌套的情况。
    \item \textbf{选择公理}：如果有很多个非空集合，那么可以从每个集合中各选出一个元素。
\end{itemize}

上述公理看似纷繁复杂，后续本书中出现的一些概念和定义也同样如此。这也是许多读者在学习数学时感到抽象和困难的重要原因之一。需要注意的是，这些公理和定义并不是凭空出现的，而是人们为了使理论体系更加严谨、完整和自洽而作出的规定。正因为要兼顾一般性与严密性，它们在形式上往往不可避免地显得繁琐。因此，在学习这些内容时，我们不必一开始就拘泥于每一个细节，而应首先把握其背后的核心思想：它想解决什么问题，刻画什么对象，以及为什么需要这样的规定。

比如上述公理中的正则公理，虽然看起来很莫名其妙，但其实它的作用是防止集合论中出现一些荒谬的情况。比如，我们可以构造一个这样的``集合''：
\[
    R=\{x\mid x\notin x\},
\]
即一个包含所有``不包含自身的集合''的集合。那么，\(R\) 是否包含自身呢？如果\( R \in R \)，根据定义则有\(R \notin R\)；反之，如果 \( R \notin R \)，根据定义则有 \(R \in R\) 。无论哪种情况都会导致矛盾，这中矛盾被称为\textbf{罗素悖论}，是集合论中最著名的悖论之一。或者，我们可以构造一个``包含所有集合的集合''\( U \)，由于 \( U \) 本身也是一个集合，因此 \( U \in U \)，又由于 \( U \) 包含所有集合，因此 \( U \notin U \)，同样导致矛盾。正则公理就是为了排除这种悖论而引入的。

好在，这种极端情况在实际应用中几乎不会遇到，因此读者不必过于担心。用朴素的观念去理解并使用集合，通常是没有问题的。但我们始终要保持警惕，意识到这些概念背后可能存在的陷阱和矛盾，并且在必要时回过头来审视它们的定义和公理，以确保我们所使用的理论体系是自洽的。

常见的集合有自然数集、整数集、有理数集、实数集和复数集，在本书中使用如下的符号来表示：
\[
    \mathbb{N},\quad \mathbb{Z},\quad \mathbb{Q},\quad \mathbb{R},\quad \mathbb{C}.
\]

集合常见的运算除了常见的交、并和差外：
\begin{align*}
    A\cap B
     & =\{x\mid x\in A\ \text{且}\ x\in B\},    \\
    A\cup B
     & =\{x\mid x\in A\ \text{或}\ x\in B\},    \\
    A\setminus B
     & =\{x\mid x\in A\ \text{且}\ x\notin B\}.
\end{align*}
其中，\(A\cap B\) 称为 \(A\) 与 \(B\) 的\textbf{交集}，\(A\cup B\) 称为二者的\textbf{并集}，\(A\setminus B\) 称为 \(A\) 关于 \(B\) 的\textbf{差集}。若 \(A\cap B=\varnothing\)，则称 \(A\) 与 \(B\) 不相交。

本书还使用到笛卡尔积。

集合本身不记录元素的顺序，但坐标需要区分不同位置。例如，平面上的点 \((1,2)\) 与 \((2,1)\) 通常不是同一个点。像 \((a,b)\) 这样带有顺序的两个元素称为\textbf{有序对}，并规定
\[
    (a,b)=(c,d)
    \quad\Longleftrightarrow\quad
    a=c\ \text{且}\ b=d.
\]

\begin{definition}[笛卡尔积]
    集合 \(A\) 与 \(B\) 的笛卡尔积定义为
    \[
        A\times B=\{(a,b)\mid a\in A,\ b\in B\}.
    \]
\end{definition}

例如，平面可以写成
\[
    \mathbb{R}^2=\mathbb{R}\times\mathbb{R}
    =\{(x,y)\mid x,y\in\mathbb{R}\}.
\]
类似地，
\[
    \mathbb{R}^n
    =\underbrace{\mathbb{R}\times\cdots\times\mathbb{R}}_{n\text{ 个}}
\]
表示所有 \(n\) 个实数组成的有序数组的集合。


在集合的基础上，映射描述了一个集合中的元素如何被送到另一个集合中。

\begin{definition}[映射]
    设 \(A,B\) 是两个集合。若按照某个确定的规则，集合 \(A\) 中的每个元素 \(x\) 都对应集合 \(B\) 中唯一的元素 \(y\)，则称这一对应关系为从 \(A\) 到 \(B\) 的映射，记作
    \[
        f:A\longrightarrow B,\qquad x\longmapsto f(x).
    \]
    其中，\(A\) 称为 \(f\) 的\textbf{定义域}，\(B\) 称为 \(f\) 的\textbf{陪域}，\(f(x)\) 称为 \(x\) 在 \(f\) 下的\textbf{像}。
\end{definition}

“函数”与“映射”在本书中含义相同。对于实数之间的对应关系，人们更常使用“函数”；当定义域或陪域是向量、矩阵等一般对象的集合时，则更常使用“映射”。

一个映射由定义域、陪域和对应规则共同确定。例如，
\[
    f:\mathbb{R}\longrightarrow\mathbb{R},
    \qquad f(x)=x^2
\]
与
\[
    g:\mathbb{R}\longrightarrow[0,+\infty),
    \qquad g(x)=x^2
\]
虽然具有相同的计算公式，但由于陪域不同，它们是两个不同的映射。一般地，两个映射相等，是指它们具有相同的定义域和陪域，并且对定义域中的每个元素都给出相同的像。


映射不仅作用于单个元素，也可以作用于集合。

\begin{definition}[像集与原像集]
    设 \(f:A\to B\)，\(S\subseteq A\)，\(T\subseteq B\)。集合 \(S\) 在 \(f\) 下的像定义为
    \[
        f(S)=\{f(x)\mid x\in S\}.
    \]
    集合 \(T\) 在 \(f\) 下的原像定义为
    \[
        f^{-1}(T)=\{x\in A\mid f(x)\in T\}.
    \]
    特别地，整个定义域的像
    \[
        f(A)=\{f(x)\mid x\in A\}
    \]
    称为 \(f\) 的\textbf{像集}。
\end{definition}
